%% CertificateGSLT: Compositional Proof Checking over Validated GSLT Presentations
%%
%% Render with:
%%   pdflatex CertificateGSLT.tex && pdflatex CertificateGSLT.tex

\documentclass[]{article}
\usepackage{url}
\usepackage[hyperindex,breaklinks]{hyperref}
\usepackage{listings}
\lstset{basicstyle=\ttfamily\footnotesize,breaklines=true,frame=single,columns=fullflexible}
\usepackage{float}
\usepackage{longtable}
\usepackage{booktabs}
\usepackage{amsfonts}
\usepackage{amssymb}
\usepackage{amsmath}
\usepackage{amsthm}

\newtheorem{theorem}{Theorem}
\newtheorem{lemma}[theorem]{Lemma}
\newtheorem{definition}[theorem]{Definition}
\newtheorem{corollary}[theorem]{Corollary}
\newtheorem{example}[theorem]{Example}
\theoremstyle{remark}
\newtheorem{remark}[theorem]{Remark}

\newcommand{\code}[1]{\texttt{#1}}
\newcommand{\codepath}[1]{\path{#1}}
\newcommand{\GSLT}{\textup{GSLT}}
\newcommand{\MeTTaIL}{\textup{MeTTaIL}}
\newcommand{\OSLF}{\textup{OSLF}}
\newcommand{\CertificateGSLT}{\textup{CertificateGSLT}}
\newcommand{\catname}[1]{\ensuremath{\mathbf{#1}}}
\newcommand{\od}[3]{\ensuremath{#2 \vdash_{#1} #3}}
\newcommand{\sem}[1]{\ensuremath{[\![#1]\!]}}

\title{\CertificateGSLT{}: Compositional Proof Checking over Validated \GSLT{} Presentations\\
\large Open Derivations, DAG Evidence, and Multisorted-Clone Semantics\\
\texttt{(DRAFT --- \today)}}
\author{Oru\v{z}i \and Zarathustra Goertzel}

\begin{document}
\maketitle

\begin{abstract}
A \CertificateGSLT{} is a validated inference presentation --- a finite family of
judgment forms and generating rules layered over a validated
\code{LanguageDef} syntax --- regarded as a finite presentation of a
proof-relevant multisorted theory.  We present a Lean~4 formalization of the
category theory and executable proof checking of such presentations.  The
development contains: a strict rule-retaining preorder of presentations whose
arrows transport derivations, checked proofs, and even the identical
chronological proof-DAG artifact; a compiled counterexample showing that
strict rule retention cannot express every legitimate translation; a calculus
of \emph{open derivations} with positional premise occurrences and a proved
simultaneous-substitution structure (both unit laws and associativity); a
provably non-thin category of derivation-valued interpretations, in which one
primitive rule may be implemented by a composite proof; set-valued models
with contravariant pullback and a semantic-transport theorem; covariant
indexed derivation families and a contravariant model family, each with a
concrete category of elements; fuel-free checkers for untrusted open proof
trees and chronological proof DAGs with premise holes, both reconstructing
\emph{exact} typed derivations; and an instantiation of a proof-relevant
multisorted clone --- the algebraic core of a cartesian multicategory --- by
the open derivations of every presentation.  A compiled sharing canary
separates compact DAG identity from expanded proof meaning: two submitted
rule nodes expand to three rule occurrences.  On top of this boundary we
construct sharing-preserving substitution of checked artifacts: checked
DAGs plug into the premise holes of a checked DAG so that acceptance is
preserved, expansion commutes with composition exactly, and submitted node
count is exactly additive while expanded cost multiplies per citation.
All results are checked with no \code{sorry} and no custom axioms; the
central theorems report exactly the standard axioms \code{propext},
\code{Classical.choice}, and \code{Quot.sound}.  The framework also contains
a proof-carrying direct trace generator; two-sided \code{Step}/\code{Steps}
adequacy against the declarative \MeTTaIL{} step relation, proved both on a
fixture and \emph{generally} over an executable admitted fragment each of
whose gate conditions carries a compiled counterexample showing the direct
translation fails without it; a versioned chronological-article wire
semantics with round-trip, canonicality, digest-identity, exact
checker-correspondence, and conservative-evolution theorems, plus a compiled
mutation battery; a genuine
fixed-core binary coproduct of rule theories; exact certificate lower bounds
with budget monotonicity; and the two orders of Belnap's four evidence
statuses.  These results are stated at their proved scope: the adequate
trace fragment is sufficient rather than maximal, the article version names
the artifact grammar rather than its presentation, and the development does
not claim a literal arithmetical-hierarchy classification, Patterson's full
constructible-duality logic, or a colimit of all languages.
\end{abstract}

\section{Introduction}
\label{sec:intro}

A validated \code{LanguageDef} finitely presents a language: sorts,
constructors, binders, equations, and rewrite rules.  Its \emph{generated
syntactic labelled transition system} (\GSLT{}) is the derived rewriting
theory --- terms, equations, steps --- and downstream tooling (parsers,
rewrite engines, behavioral type systems) is \emph{derived} from the one
presentation rather than authored separately.  The companion formalization
\code{leanOSLF.tex} develops this pipeline for operational semantics.  This
paper develops the proof-theoretic lane of the same programme: as the
rewrite lane generates a transition system from rewrite rules, the proof
lane generates \emph{derivations} from inference rules, and a
\emph{\CertificateGSLT{}} is that generated derivational theory.

Many proof formats of interest --- Metamath-style rule databases, MM0-style
verified stacks, MeTTa-hosted proof calculi, and the rule sets authored in
\GSLT{} presentations themselves --- share one shape: a finite list of
judgment forms, a finite list of schematic inference rules, and untrusted
proof artifacts that cite those rules.  The engineering temptation is to
build one checker per format.  The \CertificateGSLT{} thesis is that a single
generic checker over validated presentations suffices, and that the
\emph{relationships} between presentations (extension, refinement,
interpretation, semantic realization) form categorical structure worth
proving once, at the framework level.

Concretely, a \CertificateGSLT{} is presented by a validated inference
presentation: judgments and rules are additional fields of the same
authoritative \code{LanguageDef} that carries the syntax, so there is no
second syntax authority, and judgments are ordinary \code{Pattern} terms of
the shared syntax IR.  The generic checker consumes exactly this data.
Everything in this paper is stated and proved against that boundary.

\subsection{Contributions}

All contributions are Lean~4 theorems and executable definitions in the
\codepath{Mettapedia/GSLT/LanguageDef/} tree; \S\ref{sec:verification} gives
the verification statement.

\begin{enumerate}
\item \textbf{Strict refinement.}  A thin category of presentations whose
  arrows retain every source rule exactly.  Arrows transport typed
  derivations, checked proofs, and accepted chronological proof DAGs
  \emph{as identical artifacts}.  A compiled negative theorem shows missing
  rules make arrows uninhabited (\S\ref{sec:strict}).
\item \textbf{A compiled inadequacy counterexample.}  A two-presentation
  fixture in which the source rule $A \vdash C$ has no strict image, yet a
  derivation-valued interpretation implements it by a genuine two-node
  target proof (\S\ref{sec:counterexample}).
\item \textbf{Open derivations.}  A proof-relevant calculus of derivations
  with typed, \emph{positional} premise holes, with proved simultaneous
  substitution: two distinct unit laws, associativity, and exchange with
  positional lookup (\S\ref{sec:open}).
\item \textbf{A non-thin interpretation category.}  Interpretations map each
  admitted source rule application to a target open derivation over the same
  ordered premises.  Composition recursively translates and plugs.  The
  category laws follow from the substitution calculus, and a compiled
  witness shows a hom-type with two extensionally different arrows
  (\S\ref{sec:interpretations}).
\item \textbf{Clone semantics.}  The open derivations of every presentation
  form a proof-relevant multisorted clone --- the algebraic core of a
  cartesian multicategory: judgments are sorts, open derivations are
  operations, premise occurrences are projections, and proof plugging is
  simultaneous substitution (\S\ref{sec:clone}).
\item \textbf{Set-valued models.}  Presentations are intensional; models
  realize them by judgment-indexed carriers and rule actions.  Closed and
  open denotation, a semantic substitution lemma, contravariant model
  pullback along interpretations, and a semantic-transport naturality
  theorem are proved (\S\ref{sec:models}).
\item \textbf{Indexed structure.}  Covariant derivation and open-derivation
  functors, a contravariant model functor, and their concrete Mathlib
  categories of elements with projection functors --- the precise sense in
  which derivations are the fibers of \CertificateGSLT{}
  (\S\ref{sec:indexed}).
\item \textbf{Executable evidence.}  Fuel-free checkers for untrusted open
  proof trees and chronological open proof DAGs (premise references plus
  backward node references), with \emph{exact} reconstruction: acceptance
  yields a typed open derivation whose erasure is the submitted artifact,
  and every erasure is accepted (\S\ref{sec:checker},
  \S\ref{sec:dag}).
\item \textbf{The sharing boundary.}  Strict refinement transports the
  identical DAG artifact; a compiled canary with two submitted nodes and
  three expanded rule occurrences keeps DAG identity and cost permanently
  distinct from expanded proof meaning (\S\ref{sec:transport}).
\item \textbf{Sharing-preserving artifact substitution.}  Checked DAGs
  substitute into the premise holes of a checked DAG: the composite is
  accepted, expansion commutes with composition exactly, submitted node
  count is exactly additive while expanded cost multiplies per citation,
  and two compiled refuters reject the unshifted and uniformly wired
  variants (\S\ref{sec:dagsubstitution}).
\item \textbf{A fail-closed trace generator and an operational weld.}
  Unsupported premise kinds cannot be silently erased: a
  \code{DirectTraceLanguage} carries the premise-fragment proof required by
  \code{stepPresentation}.  A compiled counterexample rejects free
  congruence, and a nontrivial fixture proves two-sided one-step and trace
  adequacy against \MeTTaIL{}'s declarative reduction, including the
  proof-relevant diamond correspondence (\S\ref{sec:trace}).
\item \textbf{General trace adequacy over a certified fragment.}  Two-sided
  \code{Step}/\code{Steps} adequacy holds for \emph{every} admitted
  language passing an executable gate: sequentially moded, hole-skeleton,
  congruence-only, non-reflective rewrite rules with unique names.  Compiled
  counterexamples for three independent failure modes show the translation
  fails when any one of those gates is removed; the fragment is sufficient
  and is not claimed maximal (\S\ref{sec:generaladequacy}).
\item \textbf{A versioned article wire semantics.}  Closed chronological
  proof DAGs receive a canonical symbolic encoding with round-trip,
  canonicality, and digest-identity theorems; a fail-closed versioned
  checker; and an exact correspondence --- some accepted article claims a
  goal iff the goal is derivable, with linearization discharging the
  tree-versus-DAG replay obligation.  A mutation battery rejects every
  structural and semantic tamper the format must exclude, and a
  conservative-evolution law transports the identical article along strict
  rule retention, with its converse compiled as a refutation
  (\S\ref{sec:wire}).
\item \textbf{Exact finite-certificate and budget results.}  For every
  validated presentation, derivability is equivalent to existence of a raw
  proof accepted by the Boolean checker.  On the adequacy fixture, a tower
  family has accepted certificates whose rule-node count is bounded below by
  the tower height, and a checked example proves that budget exhaustion is
  not refutation (\S\ref{sec:certificates}).
\item \textbf{Fixed-core coproducts and paired evidence.}  Validated,
  identifier-disjoint rule tables over one language core form a Mathlib
  binary coproduct in the strict category.  Independent positive and negative
  certificate lanes yield Belnap's four evidence statuses, both Belnap
  orders, and status-level strong negation; a database fixture inhabits all
  four cells and separates absent ex falso from an authored ex-falso instance
  (\S\ref{sec:joins}, \S\ref{sec:four}).
\end{enumerate}

\subsection{Position in the formalization suite}

\begin{verbatim}
leanOSLF ------- categorical GSLT foundations (syntax, OSLF types)
    |
    v
CertificateGSLT ------ generic proof-checkable semantics of presentations
    |    |
    |    +-----> MeTTaIL companion (LanguageDef, engines, instances)
    |    +-----> Metamath / MM0 instantiation targets
    |
    +----------> CeTTa / Prime proof packs, MettaKernel curriculum
\end{verbatim}

\section{Validated proof presentations}
\label{sec:presentations}

The trusted boundary is the generic inference checker
(\codepath{Mettapedia/GSLT/LanguageDef/InferenceChecker.lean}).  A
\code{Presentation} is a derived view of the proof-calculus fields of the
authoritative \code{LanguageDef}: a list of judgment declarations (outer
head and arity) and a list of \code{RuleSchema} values, each with declared
metavariables, ordered premises, and one conclusion, all over the shared
\code{Pattern} syntax.

Validation is two-stage.  V1 fixes the instantiation discipline: every
metavariable has one exact occurrence depth, instantiation replaces only
declared free metavariables at that depth, collection-rest metavariables are
rejected, and binder metadata must be canonical --- no implicit lifting is
performed.  V2 adds the finite structural judgment signature: every premise
and conclusion resolves against a declared judgment head and arity, and
fixed nested constructor occurrences resolve against the underlying
\code{LanguageDef}.  A \code{ValidatedPresentation} is a presentation with
\code{isValidV2 = true}, and

\begin{lstlisting}
structure Object where
  presentation : ValidatedPresentation
\end{lstlisting}

\noindent
is the object of every category in this paper
(\codepath{Mettapedia/GSLT/LanguageDef/CertificateGSLT.lean}).

Three layers of proof data sit over an object, from trusted to untrusted:

\begin{center}
\begin{tabular}{lll}
\toprule
Layer & Lean form & Status \\
\midrule
Typed derivation & \code{Derivation P goal} & trusted, \code{Type}-valued \\
Checked proof & \code{CheckedProof P goal} & raw tree + acceptance proof \\
Raw artifact & \code{RawProof}, DAG blocks & untrusted wire data \\
\bottomrule
\end{tabular}
\end{center}

\noindent
\code{RuleApplication P i ps c} is the checker's evidence that rule instance
\code{i} admits ordered premises \code{ps} and conclusion \code{c}; it is
equivalent to successful executable instantiation
(\code{instantiateRule?\_eq\_some\_iff\_application}), which is the hinge
between the trusted and executable layers throughout.

\section{Strict refinement: the rule-retaining preorder}
\label{sec:strict}

The first category is deliberately austere.  A strict arrow retains every
source rule with the same identifier and schema:

\begin{lstlisting}
structure Morphism (source target : RuleRetaining) : Type where
  refines : RuleLookupRefines source.presentation target.presentation
\end{lstlisting}

\noindent
Strict arrows are proof-irrelevant (\code{Subsingleton}), so
\code{RuleRetaining} is a \emph{thin} category --- a preorder.  This
thinness is not an accident but a theorem-shaped statement about checking:
an arrow exists exactly when every source rule lookup succeeds identically
in the target, and then everything the checker accepted before is accepted
after.  Three transport results make that precise:

\begin{itemize}
\item \code{mapDerivation} : typed derivations transport along
  \code{refines};
\item \code{mapCheckedProof} : checked proofs transport \emph{retaining the
  exact raw tree} (\code{mapCheckedProof\_payload} is \code{rfl});
\item \code{CheckedDAG.transportMeaning} : an accepted chronological proof
  DAG re-expands to the same raw proof, which the target accepts, with a
  transported typed derivation --- the naturality square of the generic DAG
  checker.
\end{itemize}

Validated append-only presentation extensions embed as strict arrows
(\code{ofValidatedExtension}).  On the negative side, the category does not
degenerate into an indiscrete preorder:

\begin{theorem}[\code{hom\_isEmpty\_of\_missing}]
If the source presentation binds a rule identifier that the target does not,
the strict hom-type is uninhabited.
\end{theorem}

\section{Why strict retention is not enough}
\label{sec:counterexample}

Strict refinement cannot express translations that \emph{implement} a rule
rather than copy it.  The compiled fixture
(\codepath{Mettapedia/GSLT/LanguageDef/CertificateGSLTInterpretationCanary.lean})
has a source presentation with one rule $A \vdash C$ and a target
presentation with rules $A \vdash B$, $B \vdash C$, and a two-premise
sharing rule $B, B \vdash C$ --- but no rule named by the source identifier.

\begin{theorem}[\code{strict\_rule\_retaining\_interpretation\_is\_empty}]
The strict hom-type from the source to the target is uninhabited.
\end{theorem}

\begin{theorem}[\code{general\_interpretation\_exists},
\code{target\_implementation\_has\_two\_rule\_nodes}]
A derivation-valued interpretation from source to target exists, and it
implements the source rule by a target proof with exactly two primitive rule
nodes ($B\vdash C$ applied to $A \vdash B$).
\end{theorem}

\noindent
This pair of theorems makes open proof templates load-bearing: any adequate
category of proof-system translations must let a primitive rule map to a
composite proof, and composite proofs need typed premise holes.

\section{Open derivations and simultaneous substitution}
\label{sec:open}

An \emph{open derivation} $\od{P}{\Gamma}{A}$ derives judgment $A$ from an
ordered context $\Gamma$ of premise \emph{occurrences}
(\codepath{Mettapedia/GSLT/LanguageDef/CertificateGSLTOpen.lean}):

\begin{lstlisting}
inductive OpenDerivation (P : ValidatedPresentation)
    (context : List Pattern) : Pattern -> Type where
  | assumption (index : Fin context.length) :
      OpenDerivation P context (context.get index)
  | byRule (i : RuleInstance) (application : RuleApplication P i ps c)
      (children : OpenDerivationList P context ps) :
      OpenDerivation P context c
\end{lstlisting}

Premise holes are addressed by position (\code{Fin}), not by judgment:
two equal judgments at different positions are different occurrences.  A
position may be cited never, once, or repeatedly, so weakening and
contraction are present by construction rather than postulated.  This
positional discipline is exactly the cartesian structure that
\S\ref{sec:clone} names.

An \emph{environment} $\sigma : \Delta \vdash_P \Gamma$ is an ordered vector
of open derivations, one per position of $\Gamma$.  Substitution
\code{bind} plugs an environment into every hole.  The proved laws, writing
$d[\sigma]$ for \code{d.bind $\sigma$} and $\mathrm{id}_\Gamma$ for the
identity environment of premise assumptions:

\begin{align*}
x_i[\sigma] &= \sigma_i
  && \text{(\code{assumption\_bind}, \code{get\_bind})} \\
d[\mathrm{id}_\Gamma] &= d
  && \text{(\code{bind\_assumptionEnvironment})} \\
\mathrm{id}_\Gamma[\sigma] &= \sigma
  && \text{(\code{assumptionEnvironment\_bind})} \\
d[\sigma][\tau] &= d[\sigma[\tau]]
  && \text{(\code{bind\_assoc})}
\end{align*}

The two unit laws are deliberately distinct statements: the right unit says
substituting the identity environment into a derivation changes nothing; the
left unit says plugging any environment into the vector of premise
assumptions returns that environment.

Closed derivations embed into any context (\code{ofClosed}), open
derivations over the empty context close (\code{close}), and the two maps
are mutually inverse (\code{close\_ofClosed}, \code{ofClosed\_close}), so
closed proof theory is literally the empty-context fragment.

\section{The non-thin category of interpretations}
\label{sec:interpretations}

\begin{definition}[\code{Interpretation}]
An interpretation $F : P \to Q$ maps every admitted source rule application
with ordered premises $\vec{p}$ and conclusion $c$ to a target open
derivation $\od{Q}{\vec{p}}{c}$.
\end{definition}

Translation \code{mapOpen} is the evident recursion: premise occurrences map
to themselves, and a rule node maps to the interpreting template with the
translated children plugged into its holes.  The substitution-naturality law

\[
F(d[\sigma]) = F(d)[F(\sigma)]
\qquad \text{(\code{mapOpen\_bind})}
\]

\noindent
is proved by mutual induction with associativity, and the category laws of

\begin{lstlisting}
instance : Category Object where
  Hom  := Interpretation
  id   := Interpretation.id
  comp := Interpretation.comp
\end{lstlisting}

\noindent
reduce to the unit and associativity laws of \S\ref{sec:open}.  Identity
interprets each rule by itself over one distinct hole per premise;
composition interprets, then recursively reinterprets.  Strict arrows embed
by \code{ofRuleLookupRefines}, interpreting each rule by its transported
one-node template.

Unlike the strict preorder, this category is \emph{provably} not thin.  In
the fixture of \S\ref{sec:counterexample}, the source rule also has a
three-node implementation through the sharing rule
($B,B \vdash C$ applied to two copies of $A \vdash B$):

\begin{theorem}[\code{parallel\_interpretations\_distinct},
\code{interpretation\_hom\_not\_subsingleton}]
The two-node and three-node interpretations are distinct parallel arrows:
evaluating both at the canonical source rule application yields
implementations with primitive rule counts $2$ and $3$.  Consequently the
hom-type is not a subsingleton, so the interpretation category is not thin.
\end{theorem}

\begin{remark}
Proof identity, not merely provability, is the subject matter here.  A thin
quotient of this category would identify proof translations with different
costs and different sharing behavior; \S\ref{sec:transport} shows why that
identification must not be baked in.
\end{remark}

\section{The internal calculus: a proof-relevant multisorted clone}
\label{sec:clone}

The framework-level algebra of open derivations has a standard name.  A
\emph{multisorted clone} over a sort set $S$ assigns to every ordered
context $\Gamma \in \mathrm{List}\,S$ and sort $A$ a type of operations
$\mathrm{Hom}(\Gamma, A)$, with positional projections and simultaneous
substitution satisfying the projection, identity, and associativity laws
(\codepath{Mettapedia/GSLT/LanguageDef/MultiSortedClone.lean}).  Because
positions may be dropped or repeated, this is the algebraic core of a
\emph{cartesian} multicategory, as opposed to the linear discipline of a
plain multicategory.

\begin{theorem}[\code{derivationClone}]
For every \CertificateGSLT{} object, open derivations form a multisorted clone
over \code{Pattern}: sorts are ground judgment patterns, operations are open
derivations, projections are premise occurrences, and substitution is proof
plugging.  All clone laws are instances of the theorems of
\S\ref{sec:open}.
\end{theorem}

The scope of this statement is deliberately structural.  The clone supplies
the context discipline shared by every presented proof system ---
projection, weakening, contraction, cut-as-substitution.  Logical
connectives, binders, conversion rules, and side judgments are \emph{content
of the presented object theory}: they live in the rules of a particular
presentation, not in the framework.  \CertificateGSLT{} is in this sense a
logical-framework-shaped theory whose internal calculus is the substitution
structure itself.

\section{Set-valued models and semantic transport}
\label{sec:models}

A presentation is intensional: it says how derivations are \emph{generated}.
A \code{Model} realizes it extensionally
(\codepath{Mettapedia/GSLT/LanguageDef/CertificateGSLTSemantics.lean}): a carrier
type for every judgment pattern, and an action of every admitted rule
application sending realizations of the ordered premises to a realization of
the conclusion.

Closed derivations fold into any model (\code{denote}); open derivations
denote functions of a premise environment (\code{denoteOpen}).  The
semantic substitution lemma interprets plugging as composition:

\[
\sem{d[\sigma]}_\rho
  = \sem{d}_{\sem{\sigma}_\rho}
\qquad \text{(\code{denoteOpen\_bind})}
\]

Models pull back \emph{contravariantly} along interpretations: given
$F : P \to Q$ and a $Q$-model $M$, the $P$-model $F^{*}M$ keeps the carriers
and interprets each source rule by the denotation of its interpreting
template.  The functor laws \code{pullback\_id} and \code{pullback\_comp}
are proved, and semantic transport is a naturality theorem:

\begin{theorem}[\code{denote\_mapDerivation}]
$\sem{F(d)}_M = \sem{d}_{F^{*}M}$ for every
closed source derivation $d$ and target model $M$; likewise for open
derivations under any environment (\code{denoteOpen\_map}).
\end{theorem}

\noindent
Interpretations are therefore sound for arbitrary set-valued semantics:
translated proofs mean what the pulled-back model says the originals mean.

\section{Indexed structure and categories of elements}
\label{sec:indexed}

The slogan ``derivations are the fibers of \CertificateGSLT{}'' is made precise by
three functors
(\codepath{Mettapedia/GSLT/LanguageDef/CertificateGSLTIndexed.lean}):

\begin{itemize}
\item \code{derivationFunctor goal} $: \catname{Object} \to \catname{Type}$,
  covariant, acting by \code{mapDerivation};
\item \code{openDerivationFunctor context goal}, covariant, acting by
  \code{mapOpen} and retaining the ordered premise context;
\item \code{modelFunctor} $: \catname{Object}^{\mathrm{op}} \to
  \catname{Type}_{u+1}$, contravariant, acting by pullback.
\end{itemize}

Covariance of derivations and contravariance of models is the expected
variance split of a theory/model duality.  Context substitution inside one
fiber is the separate axis supplied by \code{bind}, and the two axes
commute (\code{bind\_natural}).

Each functor has its concrete Mathlib category of elements ---
\code{DerivationTotal}, \code{OpenDerivationTotal}, \code{ModelTotal} ---
with projection functors to the base.  An object of
\code{DerivationTotal goal} is a presentation together with one closed proof
of \code{goal}; a morphism is an interpretation whose action sends the
source proof to the target proof exactly.  These totals are the concrete
Grothendieck-construction endpoints over which a fully displayed
organization (over compositional syntax translations) will later sit.

\section{Executable evidence I: open proof trees}
\label{sec:checker}

Typed open derivations are the trusted calculus; wire artifacts are
untrusted (\codepath{Mettapedia/GSLT/LanguageDef/CertificateGSLTOpenChecker.lean}):

\begin{lstlisting}
inductive RawOpenProof where
  | premise (index : Nat)
  | node (ruleInstance : RuleInstance) (children : List RawOpenProof)
\end{lstlisting}

The checker \code{checkOpenRaw} is structurally recursive --- fuel-free ---
and fails closed on out-of-range premise references and failed rule
instantiations.  Its adequacy is exact in both directions:

\begin{theorem}[\code{checkOpenRaw\_erase};
\code{checkOpenRaw\_exact\_derivation}]
The erasure of every typed open derivation is accepted; and every accepted
raw tree is the erasure of some typed open derivation of the same goal in
the same context.
\end{theorem}

\begin{corollary}[\code{checkedOpenProof\_nonempty\_iff\_openDerivation}]
Checked open evidence exists if and only if a typed open derivation exists.
\end{corollary}

\noindent
``Exact'' matters: acceptance does not merely assert provability, it
reconstructs a derivation whose erasure is \emph{the submitted artifact},
so cost and shape claims about artifacts transfer to typed proofs.

\section{Executable evidence II: chronological proof DAGs}
\label{sec:dag}

Compact evidence shares subproofs
(\codepath{Mettapedia/GSLT/LanguageDef/CertificateGSLTOpenDAG.lean}).  A
chronological DAG is a list of blocks of nodes; each node carries a fresh
identifier, a rule instance, and children that either cite an authored
premise position or an \emph{earlier} derived node:

\begin{lstlisting}
inductive OpenDAGReference where
  | premise (index : Nat)
  | node (id : Nat)
\end{lstlisting}

Duplicate identifiers are rejected; the chronological environment makes
every node reference backward by construction, so no separate acyclicity
proposition is trusted.  Checking maintains a ghost expansion: each accepted
node stores the exact \code{RawOpenProof} tree it denotes.

\begin{theorem}[\code{checkOpenDAGBlocks\_exact\_derivation}]
An accepted DAG expands to an exact raw open proof accepted by the tree
checker, hence to a typed open derivation with that exact erasure.
\end{theorem}

\noindent
A compiled negative canary (\code{missing\_premise\_rejects}) shows an
out-of-range premise reference failing closed.  The closed-proof analogue of
this layer (checker \code{checkDAGBlocks}, used in \S\ref{sec:strict}) also
has an untrusted bottom-up hash-consing compiler
(\codepath{Mettapedia/GSLT/LanguageDef/InferenceProofDAGCompilation.lean});
its output is rechecked, never trusted.

\section{Exact transport and the sharing boundary}
\label{sec:transport}

Strict refinement preserves not only meaning but the artifact itself
(\codepath{Mettapedia/GSLT/LanguageDef/CertificateGSLTDAGTransport.lean}): every
successful source rule instantiation replays identically in the target, so
node, block, and root computations are unchanged.

\begin{theorem}[\code{checkOpenDAGBlocks\_true\_of\_ruleLookupRefines};
\code{CheckedOpenDAG.transport}]
Under strict rule retention, the same blocks and root identifier are
accepted by the target, and transport preserves them definitionally
(\code{transport\_rootId} and \code{transport\_blocks} are \code{rfl}).
\end{theorem}

General interpretations deliberately do \emph{not} satisfy artifact
transport: implementing one primitive node by a multi-node template changes
the artifact.  What they need instead is a genuine graph substitution, and
its premise-hole core is constructed below
(\S\ref{sec:dagsubstitution}).

The reason sharing must be tracked separately from meaning is compiled.  In
the fixture, one derived $B$ node is cited twice by the sharing rule:

\begin{theorem}[\code{shared\_open\_dag\_accepts};
\code{shared\_open\_dag\_expands\_exactly};
\code{sharing\_changes\_rule\_occurrence\_count}]
The two-node artifact is accepted, expands exactly to a tree that visibly
duplicates the shared subtree, and the counts differ: the submitted DAG has
$2$ rule nodes while the expansion has $3$ rule occurrences.
\end{theorem}

\noindent
Consequently DAG-node cost is not expanded-tree cost, and compact artifact
identity is not proof-tree identity.  The framework keeps three notions
separate by construction: expanded proof meaning (the typed derivation),
compact evidence identity (the DAG artifact), and any future cost or
provenance decoration over it.

\section{Sharing-preserving substitution of checked artifacts}
\label{sec:dagsubstitution}

Open derivations compose by \code{bind}; this section composes their
compact evidence
(\codepath{Mettapedia/GSLT/LanguageDef/CertificateGSLTDAGSubstitution.lean}, raw
layer in \codepath{CertificateGSLTRawSubstitution.lean}).  Given a checked DAG for
$\Gamma \vdash A$ and one checked DAG per premise position of $\Gamma$, all
over a common ambient context $\Delta$, the composite
(\code{CheckedOpenDAG.substitute}) lists every environment segment
\emph{once}, then the outer DAG with each premise citation rewired to the
root of its segment.  Identifiers are kept disjoint by shifting each
segment past the identifier bound of everything before it; only the outer
DAG's premise citations are rewired, while segments keep their ambient
citations.

Two compiled refuters pin this design against its natural competitors
(\codepath{CertificateGSLTDAGSubstitutionCanary.lean}): composing without the
shift collides identifiers and is rejected rather than silently merging
distinct rule nodes (\code{unshifted\_composition\_rejected}), and wiring
segment citations uniformly like outer holes makes a node cite itself
before it exists and is rejected
(\code{selfwired\_environment\_rejected}).

The surviving design satisfies the three laws that make it the artifact
presentation of clone substitution.  Write $D[\vec{E}]$ for the composite
and $\widehat{D}$ for the ghost expansion of an accepted artifact.

\begin{theorem}[\code{substituteDAGBlocks\_accepts};
\code{substituteDAGBlocks\_expands};
\code{CheckedOpenDAG.substitute\_expansion}]
The composite is accepted, and expansion commutes with composition
exactly:
$\widehat{D[\vec{E}]} = \widehat{D}\,[\widehat{E_1},\dots,\widehat{E_n}]$,
where the right-hand side is raw substitution of proof trees.
\end{theorem}

\begin{theorem}[\code{OpenDerivation.eraseOpen\_bind};
\code{CheckedOpenDAG.substitute\_erases\_bind}]
Erasure is a homomorphism from typed plugging to raw substitution, so
whenever the constituents erase typed open derivations, the composite's
expansion is the erasure of the plugged derivation: artifact substitution
presents the clone operation of \S\ref{sec:clone}.
\end{theorem}

\begin{theorem}[\code{substituteDAGBlocks\_nodeCount};
\code{RawOpenProof.substitute\_ruleCount};
\code{CheckedOpenDAG.substitute\_expansion\_ruleCount}]
Submitted node count is exactly additive,
$|D[\vec{E}]| = |D| + \sum_i |E_i|$, no matter how often a position is
cited; while the expanded tree pays each segment once per citation:
$\mathrm{rc}(\widehat{D[\vec{E}]}) = \mathrm{rc}(\widehat{D}) +
\sum_i c_i \cdot \mathrm{rc}(\widehat{E_i})$, with $c_i$ the citation count
of position $i$ in $\widehat{D}$.
\end{theorem}

\noindent
On the two-node fixture the packaged operation reproduces the boundary of
\S\ref{sec:transport} from the general theorems alone: compact cost $2$,
expanded cost $3$ (\code{packaged\_composition\_counts}).

Two honest boundaries remain by design.  Substitution never re-shares
\emph{across} constituents --- structurally equal nodes from different
segments keep distinct identifiers; re-hash-consing is a separate,
untrusted optimization whose output is rechecked.  And the artifact
calculus has no identity artifacts: a root must be a rule node, so identity
environments and the unit and associativity laws live at the expansion
level, where they are exactly the clone laws.

\section{Direct trace presentations and the \OSLF{} weld}
\label{sec:trace}

A computational step relation is one important proof theory, but operational
semantics does not generally grant contextual closure for free.  In
\MeTTaIL{}, a rewrite may use a substep only when its authored premise is a
\code{Premise.congruence}; quoted data, suspended continuations, and lazy
cells can therefore remain inert.  The first version of the trace generator
incorrectly emitted congruence at every constructor position.  The compiled
counterexample \code{congruence\_is\_not\_free} exhibits a language with
$a\to b$ but no step $f(a)\to f(b)$, permanently excluding that convention.

The corrected generator is fail-closed at the premise boundary.  A
\code{DirectTraceLanguage} packages a \code{LanguageDef} with evidence that
every rewrite premise belongs to the directly translated congruence
fragment.  Only this proof-carrying source is accepted by
\code{stepPresentation}; a language containing a relation-query premise is
rejected by \code{DirectTraceLanguage.ofLanguage?} rather than being
translated after silently dropping the premise.  The generated theory has
two judgments,

\[
  \mathsf{Step}(s,t), \qquad \mathsf{Steps}(s,t),
\]

one \code{Step} schema per authored rewrite, and reflexive--transitive trace
rules for \code{Steps}.  Generated-presentation validation remains a separate
gate: it checks metavariable depths, collection-rest freedom, constructor
references, and judgment signatures.

The collection-rest theorem is another exact format boundary.
\code{isValidV2\_eq\_false\_of\_collectionRest\_rule} proves that a schema
containing a bag-rest pattern cannot enter the current checker.  This does
not place bag semantics beyond \CertificateGSLT{}: one may present decomposition
and exchange as explicit inference, or extend the checker with a verified
bag-decomposition certificate.  What is ruled out is an unrecorded matching
choice inside the existing schema format.

The source-adequacy result was proved first as an \emph{instance theorem}
on a fixture containing a ground rewrite and an authored congruence rule;
\S\ref{sec:generaladequacy} lifts it to every language of an executable
fragment, and the fixture remains the concrete anchor for the modal weld
below.  For well-formed sources,

\begin{align*}
  \mathrm{Derivation}(\mathsf{Step}(s,t)) \text{ inhabited}
    &\iff s \mathrel{\code{langReduces}} t,
       &&\text{(\code{step\_adequacy})}\\
  \mathrm{Derivation}(\mathsf{Steps}(s,t)) \text{ inhabited}
    &\iff \mathrm{ReflTransGen}(\code{langReduces})(s,t).
       &&\text{(\code{steps\_adequacy})}
\end{align*}

The proof handles both directions against the actual matcher and premise
semantics.  It yields the first proof-relevant \OSLF{} square on the same
fixture:

\[
  \Diamond\varphi(s)
  \quad\Longleftrightarrow\quad
  \exists t,\; \mathrm{CheckedStep}(s,t) \land \varphi(t).
\]

This is \code{diamond\_iff\_checked\_witness}.  The box theorem presently
proves only the sound direction for checked, well-formed predecessors;
finite-predecessor evidence or an authored invariant theory is needed for a
converse.

\section{General trace adequacy over a certified fragment}
\label{sec:generaladequacy}

The fixture theorem generalizes, but the statement that does so must be
chosen carefully: two-sided adequacy is \emph{false} for arbitrary admitted
\code{DirectTraceLanguage}s, and each failure mode is now a compiled
counterexample rather than a design suspicion.  The general theorems
therefore quantify over an executable admission gate,
\code{languageDirectTraceAdequate}, chosen so that everything inside it is
proved and nothing outside it is constrained.

\begin{definition}[adequate direct-trace fragment]
A language passes the gate when its rewrite rules satisfy:
\begin{enumerate}
\item \emph{sequential modedness} --- reading premises in authored order,
  every congruence source is bound by the left-hand side and earlier
  congruence targets, and the right-hand side is bound by the end of the
  chain;
\item \emph{hole-skeleton patterns} --- rule patterns are
  constructor/binder skeletons whose metavariables occur at ambient depth
  only, with no explicit-substitution and no collection nodes;
\item every premise is an authored congruence premise
  (the \code{DirectTraceLanguage} gate);
\item the language declares no reflective presentations, and rewrite names
  are unique.
\end{enumerate}
\end{definition}

\begin{theorem}[\code{directTrace\_step\_adequacy},
    \code{directTrace\_steps\_adequacy}]
Let $L$ pass the gate and let its generated presentation be admitted.  For
checker-well-formed sources $s$,
\begin{align*}
  \mathrm{Derivation}(\mathsf{Step}(s,t)) \text{ inhabited}
    &\iff s \mathrel{\code{langReduces}_L} t,\\
  \mathrm{Derivation}(\mathsf{Steps}(s,t)) \text{ inhabited}
    &\iff \mathrm{ReflTransGen}(\code{langReduces}_L)(s,t),
\end{align*}
and declarative steps preserve checker well-formedness.  Composing with
\S\ref{sec:certificates}, one-step reduction and reachability in $L$ are
each equivalent to existence of an accepted finite certificate over the
generated presentation.
\end{theorem}

The proof aligns the two executable semantics that meet in the generated
schemas.  Soundness inverts an arbitrary accepted rule application into the
language's own matcher: a match-own-instance lemma produces matcher
bindings from checker arguments, a merge algebra for \code{mergeBindings}
carries binding agreement through the premise chain, and a fuel-monotonicity
theorem for the contextual step relation lets independently obtained premise
witnesses share one rule application.  Completeness runs the alignment in
reverse, reading checker arguments off the engine's final bindings; the
exact bridge is a proved agreement between schema instantiation and binding
application on hole-skeleton patterns.

\paragraph{Why the present gates are sufficient.}
The canary module supplies compiled counterexamples for three independent
failure modes of the direct translation, each showing that deleting that
gate alone breaks it:
\begin{itemize}
\item \code{moded\_gate\_necessary}: for the rule $\mathit{unit}\to x$ with
  $x$ unbound by the left-hand side, the generated presentation is admitted
  and certifies the ground step
  $\mathsf{Step}(\mathit{unit},\mathit{unit})$, yet the language steps
  $\mathit{unit}$ only to the free variable --- soundness fails;
\item \code{subst\_gate\_necessary}: an explicit-substitution reduct is
  evaluated by binding application but preserved structurally by schema
  instantiation, so the checked target is not the declarative target;
\item \code{bag\_gate\_necessary}: collection matching is
  multiset-commutative while schema instantiation is positional; the
  declarative step $\{\mathit{unit},\mathit{thunk}\}\to
  \mathit{pair}(\mathit{thunk},\mathit{unit})$ through a permuted match
  provably has no certificate --- completeness fails;
\item relation-query premises remain rejected both by
  \code{DirectTraceLanguage.ofLanguage?} and by the gate.
\end{itemize}

These counterexamples calibrate the claim precisely, and their scope should
be read exactly.  The non-reflective hypothesis is used essentially by the
present proof, which identifies source matching and application with the
syntactic checker operations; replacing it requires a separately proved
reflective-matcher adequacy interface, and no counterexample is offered for
it.  Rewrite-name uniqueness is a locally checkable convenience already
implied by validity of the generated presentation --- duplicate rewrite
names produce duplicate generated rule identifiers --- and is therefore not
an independent semantic gate.  What is established is that the fragment is
\emph{sufficient} and that the three exhibited axes are load-bearing for
this particular direct translation; neither maximality nor
translation-independent necessity is claimed.  Richer
translations remain open where they arrive with more machinery: a verified
rewrite-modulo certificate could readmit collection patterns, and certified
premise translations could admit relation and freshness premises
(\S\ref{sec:future}).  The formalization constrains exactly the proved
boundary and nothing else.

The positive canary exercises the features the fixture lacked: a language
whose reduct carries a binder literal, and a congruence rule discharging
\emph{two} authored premises in one application.  Its certificates are
obtained through the general theorem rather than by a hand-built fixture
proof, and a compiled negative confirms that the generated presentation
cannot certify a step the language does not take.

\section{Finite certificates, lower bounds, and budgets}
\label{sec:certificates}

At the checker boundary, a general theorem requires no operational fixture:

\begin{theorem}[\code{derivation\_nonempty\_iff\_certificate}]
For every validated presentation $P$ and goal $A$,
\[
  \mathrm{Nonempty}(\mathrm{Derivation}_P(A))
  \quad\Longleftrightarrow\quad
  \exists \pi:\code{RawProof},\; \code{checkRaw}(P,A,\pi)=\code{true}.
\]
\end{theorem}

The right side has the familiar finite-certificate shape of a
semidecidability result.  We do not call it a formal $\Sigma^0_1$
classification: that stronger statement additionally requires a natural
number coding and effective enumeration of \code{RawProof}.

The adequacy fixture supplies a quantitative family.  The $n$-fold contextual
lift $f^n(a)\to f^n(b)$ is proved uniformly by one induction
(\code{tower\_reachable}); every member has an accepted certificate; and
\code{tower\_certificate\_lower\_bound} proves that every accepted
certificate contains at least $n+1$ rule nodes.  Thus the development proves
both uniform derivability and an unbounded family of certificate-size lower
bounds without identifying those two kinds of evidence.

\code{checkWithinBudget} adds a two-valued verification status:
\code{established} or \code{undetermined}.  Established results are sound and
monotone under increasing budgets.  The theorem
\code{budget\_exhaustion\_is\_not\_refutation} gives a true step for which
every certificate at budget one is undetermined.  This is the proved budget
core; a full runtime verdict may add separately checked negative evidence and
operational interruption, but those are not constructors of this datatype.

\section{Versioned chronological articles: a wire semantics}
\label{sec:wire}

The DAG checker of \S\ref{sec:dag} fixes what evidence \emph{means}; this
section fixes how evidence \emph{travels}.  A \code{WireArticle} is the
closed-proof interchange artifact:

\begin{lstlisting}
WireArticle
|- version               -- ABI version (currently 1)
|- nodes                 -- chronological DAG nodes
|    |- node identifier
|    |- authored rule identifier
|    |- exact rule arguments, authored order
|    |- ordered child references
|- root node identifier
|- claimed target judgment
\end{lstlisting}

Chronology is load-bearing: a node may cite only strictly earlier nodes, so
cycles, self-references, and forward references fail locally, without a
global occurs check.  The carrier is a symbolic S-expression layer
(\code{WireTerm}); byte framing of that layer is deliberately below this
ABI.

Encoders are total and decoders fail closed.  Three theorems make the
encoding an identity discipline rather than a serialization convenience:
\emph{round trip} (\code{decodeArticle} of \code{encodeArticle} returns the
article); \emph{canonicality} (any term that decodes at all \emph{is} the
canonical rendering of its value, so each article has exactly one wire
form); and \emph{injectivity} (\code{encodeArticle\_inj\_iff}: equal canonical
renderings are equal articles).  This is injectivity on articles, not a
cryptographic digest theorem, and it identifies no presentation; it is the
layer on which a content hash could later be built.

\code{checkWireArticle} composes a fail-closed version gate with the
chronological DAG checker of \S\ref{sec:dag} at the empty premise context.
Its correspondence is exact in both directions:

\begin{theorem}[\code{wireArticle\_correspondence}]
For every validated presentation $P$ and goal $A$: some accepted article
stores target $A$ if and only if $\mathrm{Derivation}_P(A)$ is inhabited.
\end{theorem}

Soundness reconstructs a closed derivation from any accepted article via
exact reconstruction and the empty-context closing bridge.  Completeness is
a linearization theorem: \code{Derivation.linearize} emits a chronological
node list with sequential fresh identifiers, and an environment-invariant
induction proves the resulting article accepted with the derivation's goal
as its stored target.  Linearization also discharges the tree-versus-DAG
replay obligation: tree-shaped proofs enter the ABI as their
linearizations, and accepted articles expand back to exact trees.

A mutation battery accompanies the positives.  Over a small authored
presentation, one three-node article is accepted --- directly and through
its rendering --- and compiled negatives reject: a wrong ABI version; an
unknown rule identifier; a non-ground argument; a duplicate node
identifier; forward and self references; omitted, extra, and \emph{swapped}
children; a rule-identifier tamper; a premise reference in a closed
article; a bad root; a target mismatch; an undecodable term; and a tampered
rendering checked through decode.  The swapped-children rejection is the
ordered-bag discriminator: premise order is semantic, and the wire format
does not treat proof children as a multiset.

\paragraph{Exact-rule artifact retention.}
Articles outlive the presentation revision that produced them, so one
reuse law is proved directly at the wire boundary:

\begin{theorem}[\code{checkWireArticle\_true\_of\_ruleLookupRefines}]
If \code{RuleLookupRefines} holds from a source presentation to a target
presentation, then every article the source accepts is accepted by the
target --- the \emph{identical} article, with the same node identifiers,
rule instances, child order, and sharing.  The same transport holds through
the symbolic decode boundary
(\code{checkWireTerm\_true\_of\_ruleLookupRefines}).
\end{theorem}

The name matters.  This is retention of \emph{artifacts under exact rule
lookup}: it is not semantic conservativity, and it says nothing about source
models, source-language meaning, or equivalence of the two presentations.
Strict rule retention is exactly the hypothesis that makes it
unconditional: instantiation depends only on the looked-up schema and the
argument vector, and retention preserves every successful lookup.  The
converse fails, and is refuted rather than merely omitted:
\code{new\_article\_does\_not\_transport\_backward} exhibits an article
citing a rule added by an extension, accepted by the extension and rejected
by the base.  Evidence therefore moves forward under conservative
refinement and never backward.

This law covers exactly one mode of change.  Rule deletion, rule-identifier
reuse under altered semantics, and deliberate semantic revision are
\emph{not} covered, and should not be: a framework in which old evidence
survived an intentional revision would be unsound as a checking discipline.
Distinguishing conservative extension, semantics-preserving translation, and
intentional revision --- and attaching the right identity to each --- is the
subject of the following paragraph.

\paragraph{Local rule agreement, and article identity.}
Whole-table retention is more than an article needs and more than a large
library can supply: an article cites a handful of identifiers out of a table
that may hold tens of thousands.  The exact hypothesis is agreement on the
cited identifiers alone, and it yields a biconditional:

\begin{theorem}[\code{checkWireArticle\_iff\_articleRuleAgreement}]
If two validated presentations resolve every identifier an article cites to
the same schema, they accept exactly the same article.
\end{theorem}

Whole-table retention factors through this, and the local form is what makes
content-addressed article reuse across independently growing presentations
possible.

One further property is \emph{not} settled by the acceptance theorems above.
The checker validates every node and then requires only that the root carry
the target; it never requires nodes to be reachable from the root.  A
compiled witness confirms that an unreferenced but individually valid node
leaves acceptance unchanged.  This is not unsoundness --- the root still
reconstructs an exact derivation --- but it means the submitted-node count
and any content hash of an article are not functions of the proof, which the
compact-cost ledger of \S\ref{sec:transport} depends on.  Reachability, its
executable check, and the strengthened checker \code{checkRootedArticle} are
therefore given, with soundness inherited unchanged; a compiled canary shows
the strengthened checker keeps the intended article and rejects the padded
one.  Making rootedness an ABI \emph{requirement} additionally owes the
proof that linearised articles are always rooted, which is what would carry
the completeness half across; until then it is an available strengthening
rather than part of the format.

\paragraph{What the ABI does not yet fix.}
The version field identifies the artifact grammar and checking semantics
only.  The presentation is supplied externally to \code{checkWireArticle};
the article does not record which semantic presentation it belongs to.
Before persistent or cross-process use, an article should be enveloped by a
semantic identity --- presentation digest, language identity, and policy
version alongside the ABI version --- so that a structurally valid article
cannot be interpreted against the wrong theory.  Likewise,
\code{articleOfDerivation} is a straightforward linearization that does not
compress repeated subproofs; sharing-optimizing encoders remain untrusted
producers, free to emit any article the checker accepts.

\section{Fixed-core coproducts of rule theories}
\label{sec:joins}

Let $P_L$ and $P_R$ be validated rule tables over one syntax-and-judgment
core, with disjoint identifiers, and suppose their append $P_{L\mathbin{+}R}$
is also valid.  Exact lookup gives strict arrows

\[
       P_L \longrightarrow P_{L\mathbin{+}R}
       \longleftarrow P_R.
\]

Theorems \code{join\_refines\_left},
\code{join\_refines\_right}, and
\code{join\_refines\_of\_both} establish the preorder universal property.
The construction is additionally packaged against Mathlib's categorical
interface:

\begin{theorem}[\code{joinCofanIsColimit}]
The displayed binary cofan is a colimit in the strict rule-retaining
category; equivalently, the append is a binary coproduct of these two
objects.
\end{theorem}

This is a genuine finite colimit theorem, but its scope is important.  It is
not yet a coproduct across different syntax translations, a universal
interpreter, or a colimit of all computational systems.  Those statements
require the syntax-translating interpretation layer and a specified diagram.

\section{Paired evidence and Belnap's four statuses}
\label{sec:four}

The checker is neutral about negation.  Nevertheless, two independently
authored goals and two offered certificates form the evidence-status carrier
used by constructible-duality, strong-negation, and paraconsistent systems:

\[
\begin{array}{c|cc}
 & \text{negative absent} & \text{negative accepted}\\ \hline
\text{positive absent}   & \mathsf{undetermined} & \mathsf{refuted}\\
\text{positive accepted} & \mathsf{established}  & \mathsf{conflicted}
\end{array}
\]

The equivalence \code{fourEquivLanes} proves that this table is exactly
\code{Bool} $\times$ \code{Bool}, not merely a four-constructor analogy.

\code{FourVerdict.knowledgeLE} is Belnap's approximation order, with
undetermined below and conflicted above.  \code{FourVerdict.truthLE} is the
second order, with refuted below and established above.  Both are proved
reflexive, transitive, and antisymmetric.  Status-level strong negation swaps
established and refuted, fixes conflicted and undetermined, is involutive,
preserves the knowledge order, and reverses the truth order.  The operation
\code{pairedVerdict} checks offered certificates and its soundness theorems
reconstruct the positive derivation, negative derivation, or both, according
to the resulting status.  Conversely, erasures of supplied typed derivations
produce the corresponding established, refuted, or conflicted status.

A database canary joins a \code{Holds} table with a \code{Denied} table and
inhabits all four cells.  Conflict at one atom does not derive an unrelated
atom (\code{no\_explosion}); appending one explicit two-premise rule derives
that atom from the conflict (\code{explosion\_is\_authorable}).  The result is
not that the generic checker is intrinsically paraconsistent.  Rather, it
adds no implicit ex-falso principle: contradiction propagation is determined
by the authored rule theory.

This section formalizes an evidence carrier and its orders, not Patterson's
full logic of constructible duality~\cite{patterson1998constructible}.  A
complete account still needs judgment-level strong negation, logical
connectives, the lane-swapping interpretation, and appropriate semantics.
Likewise, graded $(n^+,n^-)$ or probabilistic evidence for PLN is a future
refinement of the Boolean lanes, not a consequence of the present theorems.

\section{Applications and instantiation targets}
\label{sec:applications}

\paragraph{\MeTTaIL{}-authored presentations.}
Because judgments and rules are fields of \code{LanguageDef}, every authored
\GSLT{} presentation that declares them is already a \CertificateGSLT{} object;
the companion \code{MeTTaIL.tex} describes the authoring pipeline.  The
framework then supplies, without further authorship: a proof calculus with
substitution, an interpretation category, model pullback, and two checked
evidence formats.

\paragraph{Metamath and MM0.}
Rule-database proof formats are the intended stress instantiations: their
presentations are large but structurally uniform, and their proof artifacts
are naturally chronological DAGs (Metamath's compressed proofs already cite
earlier steps by index).  The design intent is that both instantiate the
\emph{unchanged} generic checker, with source-adequacy obligations (that the
encoded presentation means the original database) kept in their respective
front ends.  This instantiation is designated future work
(\S\ref{sec:future}); the independent verified Metamath checker
\code{mm-lean4} in the same repository is the natural conformance partner.

\paragraph{CeTTa and Prime proof packs.}
The CeTTa runtime consumes compiled proof packs.  Compiling general
\CertificateGSLT{} interpretations and DAG evidence to that format, with a checker
correspondence theorem, is the designated systems-facing exit; the
hash-consing compiler of \S\ref{sec:dag} is the existing seam, and the
versioned article of \S\ref{sec:wire} is the interchange contract a native
replay checker consumes.

\section{Related work}
\label{sec:related}

Logical frameworks in the LF tradition~\cite{hhp1993lf} represent object
logics inside a dependently typed $\lambda$-calculus, obtaining substitution
and binding from the framework.  \CertificateGSLT{} is deliberately weaker at the
framework level: it fixes only the first-order structural core ---
positional contexts and simultaneous substitution --- and leaves binders and
conversion to the presented theory, in exchange for a small executable
checker and artifact-exact reconstruction theorems.  The little-theories
method~\cite{farmer1992little} develops mathematics as a graph of small
theories connected by interpretations; the non-thin category of
\S\ref{sec:interpretations} is a proof-relevant refinement of that picture
in which an interpretation carries the proofs it performs.  Clones and
Lawvere theories~\cite{lawvere1963functorial} are the classical algebra of
cartesian operations; multicategories originate with
Lambek~\cite{lambek1969deductive}.  Metamath~\cite{megill2019metamath}
demonstrates how far a minimal generic rule checker scales;
MM0~\cite{carneiro2020mm0} adds verified translation between proof
languages.  \CertificateGSLT{} aims the same minimalism at the \GSLT{} setting,
with the category theory of presentations proved rather than implicit.

Meseguer's rewriting logic~\cite{meseguer1992rewriting} gives computations a
proof-theoretic reading with reflexivity, congruence, replacement, and
transitivity.  The trace generator is close in form but intentionally does
not inherit unrestricted congruence: \MeTTaIL{} treats contextual stepping as
authored operational authority.  Belnap's four-valued information
states~\cite{belnap1977four} motivate the two evidence orders of
\S\ref{sec:four}.  Patterson's constructible duality
\cite{patterson1998constructible} is a substantially richer self-dual logic;
the present formalization supplies only its two-lane certificate carrier and
therefore avoids claiming an encoding or adequacy theorem for that logic.

\section{Verification status}
\label{sec:verification}

All results live in the Lean~4 tree under
\codepath{Mettapedia/GSLT/LanguageDef/}:

\begin{center}
\begin{tabular}{ll}
\toprule
Module & Content \\
\midrule
\code{CertificateGSLT.lean} & objects; strict rule-retaining category \\
\code{CertificateGSLTOpen.lean} & open derivations; substitution calculus \\
\code{CertificateGSLTInterpretation.lean} & non-thin interpretation category \\
\code{CertificateGSLTSemantics.lean} & set-valued models; pullback; naturality \\
\code{CertificateGSLTIndexed.lean} & indexed functors; categories of elements \\
\code{CertificateGSLTOpenChecker.lean} & open-tree checker; exact reconstruction \\
\code{CertificateGSLTOpenDAG.lean} & chronological open-DAG checker \\
\code{CertificateGSLTDAGTransport.lean} & exact artifact transport \\
\code{CertificateGSLTRawSubstitution.lean} & raw substitution; cost accounting \\
\code{CertificateGSLTDAGSubstitution.lean} & sharing-preserving composition \\
\code{CertificateGSLTClone.lean} & multisorted-clone instantiation \\
\code{CertificateGSLTStepPresentation.lean} & proof-carrying trace generator \\
\code{CertificateGSLTStepAdequacy.lean} & fixture adequacy; modalities; budgets \\
\code{CertificateGSLTStepAdequacyGeneral.lean} & fragment gate; general adequacy \\
\code{CertificateGSLTStepPresentability.lean} & checker-format boundary \\
\code{CertificateGSLTWireFormat.lean} & versioned article wire semantics \\
\code{CertificateGSLTRuleJoin.lean} & fixed-core binary coproduct \\
\code{CertificateGSLTConstructibleDuality.lean} & paired evidence; Belnap orders \\
\code{CertificateGSLTInterpretationCanary.lean} & compiled positives and negatives \\
\code{CertificateGSLTDAGSubstitutionCanary.lean} & composition refuters; counts \\
\code{CertificateGSLTStepPresentationCanary.lean} & congruence and premise gates \\
\code{CertificateGSLTStepAdequacyGeneralCanary.lean} & gate falsifiers; binder canary \\
\code{CertificateGSLTWireFormatCanary.lean} & article mutation battery \\
\code{CertificateGSLTConstructibleDualityCanary.lean} & four-cell database \\
\code{MultiSortedClone.lean} & clone interface \\
\code{CertificateGSLTFramework.lean} & canonical import \\
\bottomrule
\end{tabular}
\end{center}

The canonical import and the canaries build with zero warnings:

\begin{lstlisting}
lake build Mettapedia.GSLT.LanguageDef.CertificateGSLTFramework \
           Mettapedia.GSLT.LanguageDef.CertificateGSLTInterpretationCanary \
           Mettapedia.GSLT.LanguageDef.CertificateGSLTDAGSubstitutionCanary \
           Mettapedia.GSLT.LanguageDef.CertificateGSLTStepPresentationCanary \
           Mettapedia.GSLT.LanguageDef.CertificateGSLTStepAdequacy \
           Mettapedia.GSLT.LanguageDef.CertificateGSLTStepAdequacyGeneralCanary \
           Mettapedia.GSLT.LanguageDef.CertificateGSLTWireFormatCanary \
           Mettapedia.GSLT.LanguageDef.CertificateGSLTConstructibleDualityCanary
\end{lstlisting}

There are no \code{sorry}s, no \code{admit}s, no custom axioms, no wishlist
stubs, and no \code{native\_decide} anywhere in these modules.  Axiom audits
of the substitution laws, interpretation functoriality, semantic naturality,
exact checker reconstruction, DAG transport, the strict-arrow negative
theorem, the non-thinness witness, the sharing-count theorem, the general
trace-adequacy theorems, the gate falsifiers, and the article
correspondence and mutation theorems each report no axioms beyond
\code{propext}, \code{Classical.choice}, and \code{Quot.sound}.  Some
algebraic results use fewer.

\section{Limitations and future work}
\label{sec:future}

The current boundary is stated so that nothing beyond it is implicitly
assumed.

\begin{enumerate}
\item \textbf{Shared judgment representation.}  Interpretations preserve the
  common \code{Pattern} judgment IR.  This suffices for independently
  authored proof systems compiled into the shared IR, but it is not yet a
  category of arbitrary syntax translations.  The next layer adds syntax-
  and judgment-translating interpretations over the existing structural
  \code{LanguageDef} morphisms, with the same transport theorems.
\item \textbf{Interpretation of artifacts by rule templates.}  Premise-hole
  substitution of checked artifacts is constructed
  (\S\ref{sec:dagsubstitution}), with exact expansion and additive compact
  cost.  What remains is pushing a whole artifact through a
  derivation-valued \emph{interpretation} --- each rule node replaced by a
  DAG template for its image --- while preserving sharing; until then the
  framework does not claim that interpretation composition preserves
  DAG-node cost.
\item \textbf{Displayed organization.}  The categories of elements of
  \S\ref{sec:indexed} are concrete indexed endpoints.  A displayed category
  of proof presentations over compositional syntax \GSLT{}s, with reindexing
  along actual syntax translation, remains to be built.
\item \textbf{Cost and provenance decorations.}  Optional decorations over
  compact DAG evidence, kept strictly separate from semantic judgment
  identity as \S\ref{sec:transport} requires.
\item \textbf{Metamath and MM0 instantiation.}  Both through the unchanged
  generic checker, with source adequacy in the front ends and
  \code{mm-lean4} as the conformance partner.
\item \textbf{CeTTa/Prime pack compilation.}  A verified correspondence
  between framework evidence and the runtime proof-pack format.
\item \textbf{Beyond the certified trace fragment.}  General adequacy now
  covers the gated fragment of \S\ref{sec:generaladequacy}.  Relation-query,
  freshness, and quantified premises await certified premise translations;
  collection patterns await a verified rewrite-modulo certificate;
  reflective matching and configuration-level operational states (where a
  step is not term-to-term) each need their own presentations and adequacy
  arguments.  Each extension should arrive with its own
  counterexample-calibrated gate rather than a widened default.
\item \textbf{Semantic identity of articles.}  The wire version of
  \S\ref{sec:wire} names the artifact grammar only.  Persistent and
  cross-process articles need an envelope binding presentation digest,
  language identity, and policy version to the ABI version; byte framing
  and cryptographic hashing sit below the proved canonicality layer, and
  sharing-optimizing encoders remain untrusted producers.
\item \textbf{Constructible duality and graded evidence.}  The present
  Boolean status carrier needs a judgment-level negation interpretation,
  connectives and semantics before it realizes Patterson's logic.  PLN-style
  evidence additionally needs graded lane objects and proved composition
  laws rather than a Boolean relabeling.
\end{enumerate}

\begin{thebibliography}{9}

\bibitem{hhp1993lf}
R.~Harper, F.~Honsell, and G.~Plotkin.
\newblock A framework for defining logics.
\newblock \emph{Journal of the ACM}, 40(1):143--184, 1993.

\bibitem{farmer1992little}
W.~M. Farmer, J.~D. Guttman, and F.~J. Thayer.
\newblock Little theories.
\newblock In \emph{Automated Deduction --- CADE-11}, LNCS 607, pages
  567--581. Springer, 1992.

\bibitem{lawvere1963functorial}
F.~W. Lawvere.
\newblock Functorial semantics of algebraic theories.
\newblock \emph{Proceedings of the National Academy of Sciences},
  50(5):869--872, 1963.

\bibitem{lambek1969deductive}
J.~Lambek.
\newblock Deductive systems and categories II: Standard constructions and
  closed categories.
\newblock In \emph{Category Theory, Homology Theory and their Applications
  I}, LNM 86, pages 76--122. Springer, 1969.

\bibitem{megill2019metamath}
N.~D. Megill and D.~A. Wheeler.
\newblock \emph{Metamath: A Computer Language for Mathematical Proofs}.
\newblock Lulu Press, 2019.

\bibitem{carneiro2020mm0}
M.~Carneiro.
\newblock Metamath Zero: Designing a theorem prover prover.
\newblock In \emph{Intelligent Computer Mathematics (CICM 2020)}, LNCS
  12236, pages 71--88. Springer, 2020.

\bibitem{meseguer1992rewriting}
J.~Meseguer.
\newblock Conditional rewriting logic as a unified model of concurrency.
\newblock \emph{Theoretical Computer Science}, 96(1):73--155, 1992.

\bibitem{belnap1977four}
N.~D. Belnap, Jr.
\newblock A useful four-valued logic.
\newblock In J.~M. Dunn and G.~Epstein, editors, \emph{Modern Uses of
  Multiple-Valued Logic}, pages 8--37. D. Reidel, 1977.

\bibitem{patterson1998constructible}
A.~L. Patterson.
\newblock \emph{Implicit Programming and the Logic of Constructible
  Duality}.
\newblock PhD thesis, University of Illinois at Urbana--Champaign, 1998.

\end{thebibliography}

\end{document}
